\section{Full extension actions and the remaining phase boundary}\label{sec:actions}
\subsection{The discrepancy after a quadratic embedding has been chosen}
Even an equivariant quadratic embedding need not give the required equivariant isomorphism of the full finite extensions with every prescribed marking. This section isolates a conditional repair of the group-action part.

For clarity, take nondegenerate $\mu_2$-central extensions of the same quotient dimension, or restrict a larger target extension to a supplied invariant quadratic subspace. Suppose $C$ acts genuinely on the source and target extensions, fixes the center, and induces the same quotient action after identification. Let $F:E_s\to E_t$ be a center-fixing extension isomorphism inducing that quotient identification. Such an $F$ is supplied by the quadratic extension criterion, but is not initially assumed equivariant.

An automorphism of $E_t$ that fixes its center and quotient is uniquely of the form
\[
 u_\chi(x)=(-1)^{\chi(\bar x)}x,
 \qquad \chi\in\Hom_{\F_2}(W,\F_2).
 \tag{29.1}\label{eq:vertical}
\]
Indeed, the ratio between the image of $x$ and $x$ is a central sign depending only on $\bar x$, and preservation of multiplication makes it additive. Conversely every such character gives an automorphism. These vertical automorphisms form the additive $\F_2$-space $A=\Hom(W,\F_2)$.

Write $\alpha_g$ for the target action and $\beta_g=F\alpha^s_gF^{-1}$ for the transported source action. Their discrepancy is
\[
 \delta_g=\beta_g\alpha_g^{-1}\in A.
\]
The action of $C$ on $A$ is conjugation by $\alpha_g$. Using additive notation in $A$, the genuine group-action equations give
\[
 \delta_{gh}=\delta_g+g\delta_h.
 \tag{29.2}\label{eq:cocycle}
\]
Thus the discrepancy is an actual 1-cocycle, rather than an arbitrary list of mismatched operators.

\subsection{An explicit odd-order averaging repair}
\begin{proposition}[Conditional odd-action correction]\label{prop:odd-action}
Under the preceding hypotheses, suppose $|C|$ is odd. Put
\[
 u=\sum_{h\in C}\delta_h\in A.
 \tag{29.3}\label{eq:average}
\]
Replacing $F$ by $u\circ F$, with $u$ interpreted through \eqref{eq:vertical}, makes the extension isomorphism $C$-equivariant.
\end{proposition}
\begin{proof}
Sum \eqref{eq:cocycle} over $h$. Left multiplication permutes $C$, so
\[
 u=|C|\delta_g+gu=\delta_g+gu
\]
in the characteristic-two space $A$. Thus $\delta_g=u+gu$. Replacing $F$ by $uF$ changes its discrepancy to $u+\delta_g+gu$, which is zero for every $g$.
\end{proof}

The correction fixes the scalar center and the quotient. It may change a chosen lift by a central sign. If all generating lifts or an entire section are independently fixed, that modification may be forbidden. The valid conclusion is existence of an equivariant center-fixing extension isomorphism when these vertical corrections are allowed. It is not an equivariant pointwise identification of every independently prescribed lift.

The exact suite enumerates all cocycles for a $C_3$ action on a two-dimensional $\F_2$ module and checks the averaging formula. It also retains the opposite boundary. For $C_2$ acting trivially on a nonzero $\F_2$ module, a nonzero value of $\delta$ at its generator is a 1-cocycle. Every coboundary $u+gu$ is zero. No correction of this form can remove it. Odd order is therefore a substantive hypothesis, not a cosmetic simplification.

\subsection{What has not been repaired by this finite sum}
The result starts after the arithmetic host, actual finite actions, invariant quotient subspace and an extension isomorphism have been supplied. It does not construct them from the original source diagram. A native-normalizer obstruction can prevent the initial action from existing even when the later finite cocycle would be removable.

There is a further representation-level boundary. An action by automorphisms of a Heisenberg group preserves the isomorphism class of its representation with fixed nontrivial central character. Intertwining operators can therefore be chosen, but their products can carry scalar discrepancies. Passing from those projective operators to the specified honest phase operators is a separate problem. Proposition~\ref{prop:odd-action} removes a discrepancy in the action on the \emph{finite extension group}. It does not prove vanishing of every scalar projective obstruction in a representation, nor does it choose the prescribed Gauss or transport phases.

The original full character domain, its central quasicharacter, the noncompact root, and all specified section products remain part of the arithmetic contract. A finite central sign correction is acceptable only when its interaction with that contract is explicitly checked.

\subsection{A useful order of operations}
The two finite corrections should be used in the following order. First supply an actual invariant host and a quotient-equivariant symplectic embedding. Compare its pulled-back quadratic refinement with the source. When the defect vector is isotropic, use the equivariant transvection to correct the quadratic embedding. Next lift to the finite extensions and compute the resulting action cocycle. When the stipulated hypotheses permit it, use the odd-order sum to correct that action. Finally check every prescribed lift, section, and honest representation phase that these corrections were not authorized to change.

Each stage has a certificate and a failure mode. Combining them into one phrase such as ``equivariant attachment'' would hide which input is missing. In the equal-capacity case, even the first quadratic stage needs the native Arf match. In a source problem with a fixed embedding, the transvection may be disallowed. In even-order action problems, the 1-cocycle may survive. In all cases, the same arithmetic candidate must remain the subject of the construction.

\source{The source marking note leaves nontrivial stabilizer actions outside its established attachment theorem. This section adds a conditional finite group-action correction with a complete averaging proof and explicit even-order control. It is not a claim to have constructed general arithmetic stabilizers or solved the honest phase equations.}
